% §6.4 The Perfect Filter Line
%
% Cross-reference resolution status:
%
% DEFINED here:
%   \label{sec:perfect-filter-line}             -- this section
%   \label{def:perfect-filter-line}             -- the perfect filter line
%   \label{prop:fnr-zero-closure}               -- FNR-zero axis closure proposition
%   \label{prop:fpr-zero-closure}               -- FPR-zero axis closure proposition (companion)
%
% ALREADY DEFINED (chapter 5):
%   \Cref{sec:one-sided-two-sided}              -- chapter 5 §5.4
%   \Cref{def:one-sided}                        -- chapter 5 §5.4
%
% ALREADY DEFINED (chapter 6, earlier sections):
%   \Cref{sec:union-of-sets}                    -- §6.1
%   \Cref{prop:union-fpr-fnr}                   -- §6.1
%   \Cref{sec:other-set-ops}                    -- §6.2
%   \Cref{prop:intersection-fpr-fnr}            -- §6.2
%   \Cref{prop:complement-swap}                 -- §6.2
%   \Cref{prop:difference-fpr-fnr}              -- §6.2
%   \Cref{sec:cancellation}                     -- §6.3
%   \Cref{prop:complement-axis-swap}            -- §6.3
%   \Cref{prop:mixed-intersection-escape}       -- §6.3
%
% FORWARD REFS outside chapter 6 (chapter-level labels, exist via Plan 0 stubs):
%   \Cref{ch:optimal-bernoulli-sets}            -- chapter 14
%   \Cref{ch:perfect-hash-filters}              -- chapter 15

\section{The Perfect Filter Line}
\label{sec:perfect-filter-line}

\Cref{sec:one-sided-two-sided} named the perfect filter line as the
union of the two coordinate axes of the rate plane and previewed its
relevance to chapter~6. The composition formulas of
\Cref{sec:union-of-sets} and \Cref{sec:other-set-ops}, together with
the cancellation classification of \Cref{sec:cancellation}, have given
us everything required to read the line as an algebraic object: a
distinguished subset of $[0, 1]^2$ that some operations preserve, some
move within, and some leave. This section collects the closure
results into a single statement, makes the geometric structure of the
line explicit, and points forward to the two chapters that take each
half of the line as their constructive target.

\subsection*{Definition}

\begin{definition}[The perfect filter line]
\label{def:perfect-filter-line}
The \emph{perfect filter line} is the subset of the rate plane
$[0, 1]^2$ consisting of points with at least one coordinate equal to
zero. Explicitly, the line is the union of two coordinate axes:
\[
  \mathcal{L} \;=\; \mathcal{L}_{\fnrate} \,\cup\, \mathcal{L}_{\fprate},
\]
where the \emph{FNR-zero axis} is
$\mathcal{L}_{\fnrate} = \{(\fprate, 0) : 0 < \fprate \le 1\}$ and the
\emph{FPR-zero axis} is
$\mathcal{L}_{\fprate} = \{(0, \fnrate) : 0 < \fnrate \le 1\}$. An
approximate set \emph{lies on the line} when its rate point
$(\fprate, \fnrate)$ belongs to $\mathcal{L}$.
\end{definition}

The two axes meet only at the origin $(0, 0)$, which is the exact-set
point and is excluded from $\mathcal{L}$ by the strict inequalities in
the definition. An exact set is not an approximate set, in the sense
that one of its two rates is not merely zero but unmeasured: the
construction does not face the trade-offs the rate plane parameterizes.
The perfect filter line is the set of approximate sets that achieve
certainty in exactly one direction of report; its origin neighbor
achieves certainty in both, at the cost of the space savings that
motivated approximate sets in the first place.

\subsection*{Closure properties of the FNR-zero axis}

The composition results of the previous sections specialize cleanly
when one or both operands sit on $\mathcal{L}_{\fnrate}$. We collect
the four cases.

\begin{proposition}[FNR-zero axis closure]
\label{prop:fnr-zero-closure}
Let $\ASet{A}$ and $\ASet{B}$ be FNR-zero approximate sets, so
$\fnrate_A = \fnrate_B = 0$ and the operand rate points lie on
$\mathcal{L}_{\fnrate}$. Then:
\begin{enumerate}
  \item $\ASet{A} \cup \ASet{B}$ is FNR-zero, with $\fnrate_{\cup} = 0$
    and $\fprate_{\cup} = 1 - (1 - \fprate_A)(1 - \fprate_B)$.
  \item $\ASet{A} \cap \ASet{B}$ is FNR-zero, with $\fnrate_{\cap} = 0$
    and $\fprate_{\cap} = \fprate_A \fprate_B$.
  \item $\complement \ASet{A}$ is FPR-zero, with rates
    $(0, \fprate_A)$; the complement maps the FNR-zero axis to the
    FPR-zero axis pointwise.
  \item $\ASet{A} \setminus \ASet{B}$ is FNR-zero when $\ASet{B}$ is in
    addition FPR-zero (so $\fprate_B = 0$ also), giving rates
    $(0, \fnrate_B \cdot 0 \,\textrm{terms vanish})$; otherwise the
    difference generically leaves the line.
\end{enumerate}
In short, $\mathcal{L}_{\fnrate}$ is closed under union and
intersection, mapped to $\mathcal{L}_{\fprate}$ by complement, and only
closed under difference under a strong condition on the second operand.
\end{proposition}

\begin{proof}
Parts (1) and (2) substitute $\fnrate_A = \fnrate_B = 0$ into
\Cref{prop:union-fpr-fnr} and \Cref{prop:intersection-fpr-fnr}. The
union FNR formula~\eqref{eq:union-fnr} of \Cref{sec:union-of-sets}
becomes $w_1 \cdot 0 \cdot (1 - \fprate_B) + w_2 \cdot 0 \cdot (1 -
\fprate_A) + w_3 \cdot 0 \cdot 0 = 0$, and the intersection FNR
formula~\eqref{eq:intersection-fnr} becomes
$1 - (1 - 0)(1 - 0) = 0$. The two FPR formulas follow without
modification because they do not involve the FNRs of the operands.

Part (3) is the special case of \Cref{prop:complement-swap} with
$\fnrate_A = 0$. The swap sends $(\fprate_A, 0)$ to $(0, \fprate_A)$,
which is a point on $\mathcal{L}_{\fprate}$.

Part (4) substitutes the operand rates into
\Cref{prop:difference-fpr-fnr}. With $\fnrate_A = \fnrate_B = 0$ the
difference FNR formula becomes $\fnrate_{\setminus} = 1 - (1 - 0)(1 -
\fprate_B) = \fprate_B$, which is positive whenever $\fprate_B > 0$,
so the difference leaves $\mathcal{L}_{\fnrate}$. Imposing $\fprate_B
= 0$ in addition makes $\ASet{B}$ a perfect filter on
$\mathcal{L}_{\fprate}$, the difference FNR collapses to $0$, and
$\ASet{A} \setminus \ASet{B}$ reduces to an exact-set test.
\end{proof}

Union and intersection are the two operations that close
$\mathcal{L}_{\fnrate}$ on its own; complement closes the line
$\mathcal{L}$ as a whole by exchanging the two axes; and difference
is conditionally closed, requiring the second operand to be a perfect
filter on the opposite axis. The cancellation analysis of
\Cref{sec:cancellation} already recorded the geometric reason
difference fails: it is intersection with a complemented operand, and
the intermediate complement walks the second operand across the line.

\subsection*{Closure properties of the FPR-zero axis}

The symmetric statement for $\mathcal{L}_{\fprate}$ follows from
$\mathcal{L}_{\fnrate}$ by applying the complement-swap to both
operands of any composition and unswapping the result. We record the
companion proposition without re-deriving each case.

\begin{proposition}[FPR-zero axis closure]
\label{prop:fpr-zero-closure}
Let $\ASet{A}$ and $\ASet{B}$ be FPR-zero approximate sets, so
$\fprate_A = \fprate_B = 0$. Then $\ASet{A} \cup \ASet{B}$ has
$\fprate_{\cup} = 0$, and $\ASet{A} \cap \ASet{B}$ has $\fprate_{\cap}
= 0$; both compositions are FPR-zero. The complement of an FPR-zero
set is FNR-zero, mapping $\mathcal{L}_{\fprate}$ to
$\mathcal{L}_{\fnrate}$ pointwise. The difference $\ASet{A} \setminus
\ASet{B}$ is FPR-zero only under a strong condition on the second
operand and is otherwise two-sided.
\end{proposition}

The proof is by complement-swap from \Cref{prop:fnr-zero-closure}: any
FPR-zero set is the complement of an FNR-zero set, the rate-plane
algebra is symmetric under coordinate exchange, and the conclusions
transpose accordingly.

\subsection*{The algebraic shape of the line}

The line $\mathcal{L}$ is a non-convex subset of $[0, 1]^2$: the
midpoint of any segment from a point $(\fprate, 0)$ on
$\mathcal{L}_{\fnrate}$ to a point $(0, \fnrate)$ on
$\mathcal{L}_{\fprate}$ is $(\fprate / 2, \fnrate / 2)$, which lies in
the interior of the rate square and not on either axis. The two axes
meet only at the excluded origin, so $\mathcal{L}$ is a wedge of two
open segments rather than a connected curve; it has two ``arms''
meeting at the corner the definition excises. The shape matters
because no convex combination of operations preserves the line: mixing
one operand from each axis produces a point off both, which is exactly
the conclusion \Cref{prop:mixed-intersection-escape} reached in
\Cref{sec:cancellation}.

This is the formal sense in which the perfect filter line is a
\emph{subalgebra} of the rate plane: closed under a specific subset of
the chapter~6 operations, exited by the remaining operations, and
internally structured into two halves that the complement operator
swaps. Same-axis operations compose freely without leaving the line;
cross-axis operations require care, because their composed structure
is no longer a perfect filter and the chapter~14 and chapter~15 space
guarantees do not apply.

\subsection*{Forward to chapters 14 and 15}

Two later chapters take the two halves of the line as their
constructive concern. The forward references close the picture of
why the line matters for the rest of Part~II.

\Cref{ch:optimal-bernoulli-sets} develops the Bernoulli Hash
Function, the constructive realization of an approximate set on
$\mathcal{L}_{\fnrate}$ at the lowest possible space cost. The
information-theoretic floor of chapter~8 gives the minimum number of
bits per element required to represent an FNR-zero approximate set at
false-positive rate $\fprate$; the BHF achieves the floor up to a
small additive constant, closing the $\log_2 e \approx 1.44$
multiplicative overhead the standard Bloom filter of chapter~5 leaves
above the floor on the same axis. The closure results of this section
then govern composed BHFs: union and intersection of BHFs remain on
$\mathcal{L}_{\fnrate}$ at rates the chapter~6 formulas predict, so
the entropy floor extends to the composition.

\Cref{ch:perfect-hash-filters} develops the perfect hash filter, the
constructive realization on $\mathcal{L}_{\fprate}$ for a designated
input set $A$. A perfect-hash computation at build time stores a
fingerprint per element; a query for $x \in A$ lands on its
fingerprint and is reported correctly with probability one, while a
query for $x \notin A$ may still hash to a stored fingerprint. The
construction sits on the FPR-zero axis and stays there under
same-axis union and intersection by \Cref{prop:fpr-zero-closure}.
Mixed-axis composition with a BHF, by
\Cref{prop:mixed-intersection-escape}, produces a two-sided structure
off the line and the FPR-zero certainty is lost; the designer must
respect the closure constraints to retain the constructive guarantees
of either chapter.
